% part: intuitionistic-logic
% chapter: propositions-as-types
% section: type-preservation

\documentclass[../../../include/open-logic-section]{subfiles}

\begin{document}

\olfileid{int}{pty}{tpr}

\olsection{Panglestaren Tipe}

Konversi reduksi lan permutasi sing wis dibahas uga bisa ditetepake
langsung kanggo sistem \Log{tN2i} lan \Log{tN3i}. Sistem \Log{tN3i}
netepake tipe terma $\lambd$~$N$ relatif marang konteks~$\Gamma$
minangka !!{formula}~$!A$ kanthi $\Gamma \Sequent N : !A$ nduweni
!!a{proof} ing \Log{tN3i}. Amarga aturan reduksi terma $\lambd$
nggambarake konversi reduksi lan permutasi !!{proof}s kanthi pas,
ana akibat wigati: tipe terma $\lambd$ relatif marang konteks~$\Gamma$
tetep padha sawise konversi reduksi utawa permutasi. Sipat iki diarani
\emph{panglestaren tipe}.

\begin{prop}
Yen $N$ nduweni tipe~$!A$ relatif marang $\Gamma$ lan $N \redone N'$,
$N'$ nduweni tipe~$!A$ relatif marang~$\Gamma$.
\end{prop}

\begin{proof}
  Kanthi induksi marang~$N$ lan dhuwure !!{proof}s~$\delta$ saka $\Gamma
  \Sequent N : !A$, kita bisa netepake pratelan iki: yen $M$ iku
  subterma saka~$N$, $\delta$ ngandhut sub-!!{proof}~$\delta_1$
  sing dipungkasi ing $\Gamma' \Sequent M : !B$. Yen $M$ iku redex,
  sawijining konversi bisa ditrapake marang~$\delta_1$. Kanthi
  ngetrapake reduksi marang $\delta_1$, kita ngasilake
  !!a{proof}~$\delta_1'$ saka $\Gamma' \Sequent M' : !B$.
  Kanthi mriksa konversi \Log{tN3i}, $M'$ iku asil kontraksi saka~$M$.
  Ngganti $\delta_1$ ing~$\delta$ nganggo $\delta_1'$ ngasilake
  !!a{proof}~$\delta'$ saka $\Gamma \Sequent N' : !A$, kanthi $N'$
  diasilake saka~$N$ kanthi ngganti subterma~$M$ nganggo~$M'$.

  Contone, gatekna konversi reduksi kanggo~\Intro{\land} sing
  ditututi \Elim{\land}:
  \[
  \bottomAlignProof
  \AxiomC{}
  \Deduce$\Gamma \fCenter N_1 : !B_1$
  \AxiomC{}
  \Deduce$\Gamma \fCenter N_2 : !B_2$
  \RightLabel{\Intro{\land}}
  \BinaryInf$\Gamma \fCenter \pair{N_1}{N_2} : !B_1 \land !B_2$
  \RightLabel{\Elim{\land}[i]}
  \UnaryInf$\Gamma \fCenter \proj{i}{\pair{N_1}{N_2}} : !B_i$
  \DisplayProof
  \quad\redone\quad
  \bottomAlignProof
  \AxiomC{}
  \Deduce$\Gamma \fCenter N_i : !B_i$
  \DisplayProof
  \]
  Kita ndeleng yen terma bukti saka !!{proof} ing tengen, yaiku
  $\delta_1'$, iku $N_i$. Iki persis asil ngetrapake konversi
  $\beta$ marang terma bukti $\proj{i}{\pair{N_1}{N_2}}$ saka
  !!{proof} ing kiwa~($\delta_1$).

  Utawa gatekna \Intro{\lif} sing ditututi inferensi \Elim{\lif},
  lan konversi reduksi sing ditrapake marang pasangan iku:
  \begin{multline*}
  \bottomAlignProof
  \AxiomC{$\delta_2$}
  \Deduce$x: !B_1, \Gamma \fCenter N : !B_2$
  \RightLabel{\Intro{\lif}}
  \UnaryInf$\Gamma \fCenter \lambd[\typeof{x}{!B_1}][N] : !B_1 \lif !B_2$
  \AxiomC{}
  \RightLabel{$\delta_3$}
  \Deduce$\Gamma \fCenter M : !B_1$
  \RightLabel{\Elim{\lif}}
  \BinaryInf$\Gamma \fCenter (\lambd[\typeof{x}{!B_1}][N] M) : !B_2$
  \DisplayProof
  \quad\redone\\
  \bottomAlignProof
  \AxiomC{}
  \RightLabel{$\delta_3$}
  \Deduce$\Gamma \fCenter M : !B_1$
  \RightLabel{$\Subst{\delta_2}{\delta_3}{x:!B_1}$}
  \Deduce$\Gamma \fCenter \Subst{N}{M}{x} : !B_2$
  \DisplayProof
  \end{multline*}
  Terma bukti saka !!{proof} ing tengen maneh persis asil reduksi
  $\beta$ marang terma bukti saka !!{proof} ing kiwa. Ganti jeneng
  variabel kaiket dhisik supaya variabel sing diuculake anyar tumrap
  konteks njaba, lan pangiket njero anyar tumrap terma pangganti.
  Induksi marang dhuwure~$\delta_2$ nuduhake yen saben aksioma
  $\Gamma' \Sequent x:!B_1$ ing~$\delta_2$ diganti nganggo
  !!{proof}~$\delta_3$ saka $\Gamma \Sequent M:!B_1$, sing diwenehi
  pelemahan supaya cocog karo konteks lokal sing isih ana sawise
  variabel sing diuculake dibusak. Aksioma variabel liyane tetep,
  lan saben aturan introduksi utawa eliminasi njaga tipe asile;
  syarat variabel anyar nyegah pangcekelan ing abstraksi lan cabang
  kasus. Iki ngasilake !!a{proof} $\Subst{\delta_2}{\delta_3}{x:!B_1}$
  saka $\Gamma \Sequent \Subst{N}{M}{x} : !B_2$.
\end{proof}

Korespondhensi raket antarane !!{proof}s lan terma $\lambd$ nduweni
akibat liyane. Yen $N$ iku terma $\lambd$ mawa tipe $!A$ relatif
marang konteks~$\Gamma$ lan ngandhut redex~$M$ (yaiku durung awujud
normal), bukti \Log{tN3i}~$\delta$ sing cocog saka $\Gamma \Sequent N : !A$
ngandhut sub-!!{proof} (sing dipungkasi ing $\Gamma' \Sequent M: !B$)
sing bisa diwenehi konversi. Yen $N \redone N'$ kanthi ngganti $M$
nganggo asil kontraksine, asil sing cocog saka ngetrapake reduksi
marang~$\delta$ iku !!a{proof} saka $\Gamma \Sequent N' : !A$.
Mula saben terma $\lambd$ sing tipene bener lan durung awujud normal
direduksi dadi terma liyane kanthi konversi reduksi utawa permutasi.
Sipat iki diarani \emph{kemajuan}.

Kemajuan lan panglestaren tipe njamin yen reduksi terma $\lambd$
ora tau ``mandheg tanpa langkah'': yen sawijining terma tipene bener
lan durung awujud normal, terma iku direduksi dadi terma liyane sing
tipene bener (lan nduweni tipe sing padha).

\end{document}
